\subsection{Ordinary algorithms over arithmetic strict henselizations}
\label{C:strict-henselization-subsection}

Fix a computable algebraic closure
$\Omega=\overline{\mathbf F}_p$ as a coherent union of finite
fields $\mathbf F_{p^{s!}}$.  Finite-field tables and embeddings
are found by finite search; elements have finite-stage names,
and equality and field arithmetic take place in a common stage.
The prime $p$ is odd.  A finite residue field is embedded in
this model once and for all, by a computable finite-field map.
This is the only choice of geometric residue point used below.

\begin{lemma}[Proved: exact names and full kernels for strict henselizations]
\label{C:arithmetic-strict-henselization-access}
Let $S=\mathbf Z[z_1,\ldots,z_s]/I$ be a supplied finite
presentation, and let $S\twoheadrightarrow\kappa$ be a
supplied finite residue field of characteristic $p>2$.
Put $\mathfrak m=\ker(S\to\kappa)$ and
$R=(S_{\mathfrak m})^{\rm sh}$ for the chosen
$\kappa\hookrightarrow\Omega$.
Then $R$ has finite names with computable exact arithmetic,
equality, unit inversion and residue in $\Omega$.
Any requested element of $\Omega$ has a computable lift to
$R$.  Finite generators of the full kernel of every finite
matrix over $R$ are computable by an ordinary terminating
algorithm.  These conclusions apply to singular,
nonreduced and mixed-characteristic $S_{\mathfrak m}$.
\end{lemma}
\begin{proof}
The ring $R$ is Noetherian and strictly henselian, with
residue field $\Omega$
\cite[Section 10.155, Lemma 10.155.2; Tag 06LJ,
Lemma 15.46.3]{C-StacksStrictEffective}.
We give names which retain the chosen geometric point.
Choose $b\in S$ with nonzero residue, a monic
$f\in S_b[U]$, and $a\in\Omega$ satisfying
$\bar f(a)=0$ and $\bar f'(a)\ne0$.
Choose $h\in S_b[U]$ with $\bar h(a)\ne0$, and set
\[
 T=S_b[U,V]/\bigl(f(U),\,Vf'(U)h(U)-1\bigr),
 \qquad \rho(U)=a,\qquad
 \rho(V)=(\bar f'(a)\bar h(a))^{-1}.                    \tag{C.3.SH1}
\]
Here $\rho:T\to\Omega$ extends the fixed map on $S$.
It is a geometric point, and is not asserted to be
surjective onto the infinite field $\Omega$.
Its image is a finite subfield of $\Omega$: it is
contained in a finite field containing the images
of the finitely many generators, and is a finite
subring of a field.  Thus
$\mathfrak q=\ker\rho$ is maximal, with its specified
embedding $\kappa(\mathfrak q)\hookrightarrow\Omega$.
All displayed conditions are decidable in finite
stages of $\Omega$.

A raw name is this finite presentation together
with a polynomial element of $T$.
The variable inverting $b$ is included explicitly,
so $T$ is a finitely presented integer algebra.
For intermediate computations allow finite tensor
products of such neighborhoods over $S$, using the
product geometric point in $\Omega$.
They are again finitely presented \'etale
neighborhoods, with explicitly presented rings.
Every element of $R$ has a raw name: strict
henselization is the filtered colimit of \'etale
neighborhoods with specified residue embeddings
in $\Omega$, and near the specified point they
are standard \'etale
\cite[Lemmas 10.155.2 and 10.155.11; Tag 00UE,
Proposition 10.144.4]{C-StacksStrictEffective}.
The finitely many coefficients over $S_{\mathfrak m}$
lie in one $S_b$, giving precisely (C.3.SH1).

With the induced residue embedding there is a
canonical pointed identification
\[
                      (T_{\mathfrak q})^{\rm sh}=R.       \tag{C.3.SH2}
\]
Indeed the unique lift of the geometric point of
$T$ to the strictly henselian ring $R$, and the
universal properties of the two strict henselizations,
give maps in both directions, inducing the identity
on $\Omega$.  Their compositions are identities by
uniqueness of the geometric \'etale lift and of
strict henselization
\cite[Lemmas 10.155.9--10.155.11]{C-StacksStrictEffective}.
Equivalently the pointed neighborhoods over $T$
are cofinal: tensoring any pointed $S$-neighborhood
with $T$ gives one over $T$ with the product
geometric point.  Consequently $T\to R$ is flat
and $T_{\mathfrak q}\to R$ is faithfully flat.

Use Lemma~\ref{C:integer-algebra-linear-algorithms}
to compute annihilators in $T$.  For $z\in T$,
\[
 z=0\text{ in }R
 \quad\Longleftrightarrow\quad
 \operatorname{Ann}_T(z)\not\subseteq\mathfrak q.        \tag{C.3.SH3}
\]
The right side holds exactly when at least one
computed annihilator generator has nonzero image
under $\rho$.  This is a finite exact test in
$\Omega$.  For two names use their tensor
neighborhood and test their difference.
Arithmetic is performed in that neighborhood;
a unit is characterized by nonzero residue,
and its inverse is named by adjoining an inverse
variable.  Thus this construction gives equality
in the actual strict henselization, not just a
semidecision from expansions.

The raw finite names have decidable validity.
Retaining a code precisely when no smaller valid
code has the same value gives canonical finite
names as a decidable subset of $\mathbf N$.
An internal tensor or localization name is returned
to the raw list by an equality search; coverage
proves termination.
To lift $c\in\Omega$, compute its minimal polynomial
over $\mathbf F_p$ inside a finite field containing
it, lift that monic polynomial to $\mathbf Z[U]$,
and use (C.3.SH1) with the simple geometric root
$a=c$ and $h=1$.  Its generator gives the requested
lift.  This construction does not assert the
existence of a coefficient field inside the
mixed-characteristic ring $R$.

Finally put the entries of a finite matrix $A$
in one tensor neighborhood $T$, obtaining $A_T$.
Compute generators of $\ker A_T$ by the integer
polynomial algorithms.  Flatness in (C.3.SH2)
gives
\[
                       \ker A=(\ker A_T)\otimes_T R.
\]
The computed generators therefore generate the
full actual kernel.  All computations have finite
names and the asserted termination proofs.
\end{proof}

\begin{theorem}[Proved: complete ordinary classification over arithmetic strict henselizations]
\label{C:arithmetic-strict-henselization-WG}
For $R$ as in
Lemma~\ref{C:arithmetic-strict-henselization-access},
all indecomposable finite projective rational
representations of split $O_m$ have a computable
literal canonical index set $\mathscr I_R\subseteq\mathbf N$.
Isomorphism, indecomposability, duality, invariant
form sign, equivariant isometry and equality in
the split group are decidable by ordinary
terminating algorithms.

Partition the self-dual indices into symmetric
and alternating types $\mathscr I_{\rm s}$ and
$\mathscr I_{\rm a}$, and choose the smaller
index in each non-self-dual pair to obtain
$\mathscr I_{\rm h}$.
Then actual computed generators give
\[
 WG_{\rm sp}(BO_{m,R})\simeq
       \bigoplus_{i\in\mathscr I_{\rm s}\cup
                         \mathscr I_{\rm a}\cup
                         \mathscr I_{\rm h}}\mathbf Z e_i.         \tag{C.3.SH4}
\]
The generator is an invariant symmetric form on
$M_i$ in symmetric type and $H(M_i)$ in alternating
or paired type.  All products and specified
exterior-power coefficients in (C.3.SH4) are
computed in finite time.
There are no determinant bits, and the forgetful
map to the split representation group is injective.
No polynomial-system, splitting-field, or
infinite-branch oracle is assumed.
\end{theorem}
\begin{proof}
The ring is Noetherian henselian local and
$2$ is a unit.  Its finite projectives are free.
The previous lemma computes every finite Hom
and invariant-form kernel in the coefficient
modules $C_d(R)$, and computes residues of their
generators in $\Omega$.
We explain the replacements for the finite
residue searches in
Theorem~\ref{C:effective-henselian-DVR-WG}.

\emph{Deciding invertible linear solutions.}
Let $H_1,\ldots,H_s$ generate a module of
$n$-square matrices defined by the required
linear equations over $R$.
The residues of actual solutions form the
$\Omega$-span of the $\bar H_j$, since every
scalar in $\Omega$ has a lift to $R$.
The determinant
\[
             P(c_1,\ldots,c_s)
                =\det\Bigl(\sum_j c_j\bar H_j\Bigr)
\]
has degree at most $n$ in each variable.
Choose $n+1$ distinct elements of $\Omega$,
found in one sufficiently large finite field,
and test $P$ on their finite Cartesian grid.
A polynomial of degree at most $n$ in each
variable is zero if and only if it vanishes
on this grid, by repeated one-variable interpolation.
Thus a successful test exists exactly when
there is an invertible actual solution.
Lift its finitely many scalars by the previous
lemma; the resulting actual linear combination
has unit determinant.  This decides and constructs
module isomorphisms and nonsingular reference
forms of either sign.  The rank-zero case is
handled directly.

\emph{Computing residual endomorphism algebras.}
For a module $M$, compute generators of
$E=\operatorname{End}(M)$ and their finite-dimensional
image $B$ in the residue matrix algebra over
$\Omega$.  Linear algebra over $\Omega$ computes
a basis, its multiplication constants and its
identity.  Choose a finite subfield $k_0\subset\Omega$
containing all entries of that basis and all
these structure constants.  Their $k_0$-span
is a finite algebra $B_0$ with
$B=B_0\otimes_{k_0}\Omega$.
Compute $J_0=\operatorname{rad}B_0$ by finite
enumeration and the Jacobson unit criterion.

We justify that this computes the radical after
extension as well:
\[
       \operatorname{rad}B=J_0\otimes_{k_0}\Omega.          \tag{C.3.SH5}
\]
The ideal on the right is nilpotent, hence
contained in the radical.  The finite semisimple
algebra $B_0/J_0$ is a product of matrix algebras
over finite division rings, which are finite
fields.  These fields are separable over $k_0$,
so extension to $\Omega$ splits each into a
product of copies of $\Omega$.  The resulting
quotient is semisimple, proving the reverse
inclusion in (C.3.SH5).  This uses only the
finite-field algebra facts already used in
Theorem~\ref{C:Om-finite-field-algorithm}.
The radical generators and quotient dimensions
are therefore computed by finite-field linear
algebra; a splitting oracle has not been used.

The kernel of $E\to B$ is contained in
$\operatorname{rad}E$, since an endomorphism
zero modulo the maximal ideal has $1-uh$
invertible for every $u\in E$ by its determinant.
Thus
$E/\operatorname{rad}E=B/\operatorname{rad}B$.
By Lemma~\ref{C:henselian-finite-algebra-lifting},
$E$ is semiperfect.  For a positive-rank $M$,
indecomposability is consequently equivalent
to $B/\operatorname{rad}B$ having dimension one
over $\Omega$.
Indeed a finite-dimensional division algebra
over the algebraically closed field $\Omega$
is $\Omega$ itself: each element satisfies a
polynomial which splits into central linear
factors, and a division ring has no zero divisors.
A semisimple algebra of larger dimension is
not a division algebra and has a nontrivial
idempotent.

In that larger-dimension case a nontrivial
idempotent in $B$ can be found effectively.
Enumerate larger finite subfields $k_1$ of
$\Omega$ containing $k_0$ and all elements of
$B_0\otimes_{k_0}k_1$, testing the idempotent
equation in each finite algebra.
The known nontrivial idempotent in $B$ has
finitely many coefficients, so one of these
finite searches finds it.  This search is
entered only after its existence has been
proved by the quotient dimension test.
Lift it to an actual idempotent of $E$ by
enumerating exact endomorphisms with that
residue; the henselian lifting lemma proves
termination.  Split its free image and
complement by residue columns with a unit
minor and repeat.  Positive ranks decrease,
so a full Krull--Schmidt decomposition is
computed after finitely many splits.

\emph{Indices, forms and coordinates.}
Use finite raw ring names and coaction coefficients
to code all representations by natural numbers.
Coaction identities, minimal coefficient degree
and indecomposability are decidable by the
preceding procedures.
Keep an indecomposable code precisely when no
smaller valid code is isomorphic to it.
There are finitely many predecessors and their
isomorphism tests terminate.  Every class has
one least code, which is its literal canonical
index.  Applying the antipode and the same
procedure computes duality.

For a self-dual indecomposable, the local-ring
argument applied to $g+g^{\mathsf t}$ and
$g-g^{\mathsf t}$ supplies at least one sign
of nonsingular invariant form.  The grid test
finds such a form and decides its sign.
Its residue endomorphism division algebra is
$\Omega$, and its adjoint fixes every residue
scalar: the adjoint fixes scalar multiplication
by every element of $R$, which surjects onto
$\Omega$.  Thus there is no nontrivial residue
involution.  Both signs cannot occur on the
same indecomposable: their ratio would be an
invertible skew-adjoint endomorphism, whose
residue scalar would satisfy $c=-c$, impossible
for a unit when $p>2$.

Hermitian radical lifting from
Lemma~\ref{C:henselian-finite-algebra-lifting}
now reduces each self-dual block to forms over
$\Omega$ with trivial involution.
Symmetric forms over an algebraically closed
field are classified by dimension, and alternating
forms by half their even dimension.  A non-self-dual
pair has the usual hyperbolic multiplicity.
This proves (C.3.SH4), with the stated actual
generators and no torsion bits.
It also proves that two actual symmetric forms
are isometric if and only if their underlying
representations are isomorphic.
Their coordinates are recovered from the
Krull--Schmidt multiplicities: symmetric
multiplicity itself, half the alternating
multiplicity, and the common multiplicity in
a dual pair.  If an actual isometry is requested,
enumerate exact matrices satisfying its equations
after this positive decision; existence proves
termination.

The split representation group is free on the
indecomposable indices.  The forgetful images of
the generators in (C.3.SH4) are $[M_i]$ in
symmetric type, $2[M_i]$ in alternating type,
and $[M_i]+[M_i^*]$ in paired type.
They are independent, proving the asserted
injectivity.  Finally tensor coefficient
matrices and exterior minors compute the actual
operation outputs before decomposing them.
The preceding ordinary algorithms recover their
coordinates.  Bilinearity and finite expansion
of $\lambda_t(x-y)=\lambda_t(x)\lambda_t(y)^{-1}$
handle every requested virtual exterior coefficient.
This completes the classification and termination
proof without a residual form or representation
isomorphism quotient.
\end{proof}
